\section{Algebraic closure}
\label{section-algebraic-closure}

\noindent
The ``fundamental theorem of algebra'' states that $\mathbf{C}$ is
algebraically closed. A beautiful proof of this result uses
Liouville's theorem in complex analysis, we shall give another
proof (see Lemma \ref{lemma-C-algebraically-closed}).

\begin{definition}
\label{definition-algebraically-closed}
A field $F$ is said to be {\it algebraically closed} if every algebraic
extension $E/F$ is trivial, i.e., $E = F$.
\end{definition}

\noindent
This may not be the definition in every text. Here is the lemma comparing
it with the other one.

\begin{lemma}
\label{lemma-algebraically-closed}
Let $F$ be a field. The following are equivalent
\begin{enumerate}
\item $F$ is algebraically closed,
\item every irreducible polynomial over $F$ is linear,
\item every nonconstant polynomial over $F$ has a root,
\item every nonconstant polynomial over $F$ is a product of linear factors.
\end{enumerate}
\end{lemma}

\begin{proof}
If $F$ is algebraically closed, then every irreducible polynomial is linear.
Namely, if there exists an irreducible polynomial of degree $> 1$, then
this generates a nontrivial finite (hence algebraic) field extension, see
Example \ref{example-degree-simple-algebraic-extension}.
Thus (1) implies (2). If every irreducible polynomial
is linear, then every irreducible polynomial has a root, whence every
nonconstant polynomial has a root. Thus (2) implies (3).

\medskip\noindent
Assume every nonconstant polynomial has a root. Let $P \in F[x]$
be nonconstant. If $P(\alpha) = 0$ with $\alpha \in F$, then we see
that $P = (x - \alpha)Q$ for some $Q \in F[x]$ (by division with remainder).
Thus we can argue by induction on the degree that any nonconstant
polynomial can be written as a product $c \prod (x - \alpha_i)$.

\medskip\noindent
Finally, suppose that every nonconstant polynomial over $F$ is a product of
linear factors. Let $E/F$ be an algebraic extension. Then all the simple
subextensions $F(\alpha)/F$ of $E$ are necessarily trivial (because the
only irreducible polynomials are linear by assumption). Thus $E = F$.
We see that (4) implies (1) and we are done.
\end{proof}

\noindent
Now we want to define a ``universal'' algebraic extension of a field.
Actually, we should be careful: the algebraic closure is {\it not} a
universal object. That is, the algebraic closure is not unique up to
{\it unique} isomorphism: it is only unique up to isomorphism. But still,
it will be very handy, if not functorial.

\begin{definition}
\label{definition-algebraic-closure}
Let $F$ be a field. An {\it algebraic closure} of $F$ is a field
$\overline{F}$ containing $F$ such that:
\begin{enumerate}
\item $\overline{F}$ is algebraic over $F$.
\item $\overline{F}$ is algebraically closed.
\end{enumerate}
\end{definition}

\noindent
If $F$ is algebraically closed, then $F$ is its own algebraic closure.
We now prove the basic existence result.

\begin{theorem}
\label{theorem-existence-algebraic-closure}
Every field has an algebraic closure.
\end{theorem}

\noindent
The proof will mostly be a red herring to the rest of the chapter. However, we
will want to know that it is {\it possible} to embed a field inside an
algebraically closed field, and we will often assume it done.

\begin{proof}
Let $F$ be a field. By Lemma \ref{lemma-size-algebraic-extension} the
cardinality of an algebraic extension of $F$ is bounded by
$\max(\aleph_0, |F|)$. Choose a set $S$ containing $F$ with
$|S| > \max(\aleph_0, |F|)$. Let's consider triples
$(E, \sigma_E, \mu_E)$ where
\begin{enumerate}
\item $E$ is a set with $F \subset E \subset S$, and
\item $\sigma_E : E \times E \to E$ and $\mu_E : E \times E \to E$
are maps of sets such that $(E, \sigma_E, \mu_E)$ defines the structure
of a field extension of $F$ (in particular $\sigma_E(a, b) = a +_F b$
for $a, b \in F$ and similarly for $\mu_E$), and
\item $E/F$ is an algebraic field extension.
\end{enumerate}
The collection of all triples $(E, \sigma_E, \mu_E)$ forms a set $I$.
For $i \in I$ we will denote $E_i = (E_i, \sigma_i, \mu_i)$ the
corresponding field extension to $F$. We define a partial ordering on
$I$ by declaring $i \leq i'$ if and only if $E_i \subset E_{i'}$
(this makes sense as $E_i$ and $E_{i'}$ are subsets of the same set $S$)
and we have $\sigma_i = \sigma_{i'}|_{E_i \times E_i}$ and
$\mu_i = \mu_{i'}|_{E_i \times E_i}$, in other words, $E_{i'}$ is a field
extension of $E_i$.

\medskip\noindent
Let $T \subset I$ be a totally ordered subset. Then it is clear that
$E_T = \bigcup_{i \in T} E_i$ with induced maps $\sigma_T = \bigcup \sigma_i$
and $\mu_T = \bigcup \mu_i$ is another element of $I$. In other words
every totally order subset of $I$ has a upper bound in $I$. By Zorn's lemma
there exists a maximal element $(E, \sigma_E, \mu_E)$ in $I$. We claim that
$E$ is an algebraic closure. Since by definition of $I$ the extension
$E/F$ is algebraic, it suffices to show that $E$ is algebraically closed.

\medskip\noindent
To see this we argue by contradiction. Namely, suppose that $E$ is not
algebraically closed. Then there exists an irreducible polynomial
$P$ over $E$ of degree $> 1$, see Lemma \ref{lemma-algebraically-closed}.
By Lemma \ref{lemma-finite-is-algebraic} we obtain a nontrivial finite
extension $E' = E[x]/(P)$. Observe that $E'/F$ is algebraic by
Lemma \ref{lemma-algebraic-permanence}.
Thus the cardinality of $E'$ is $\leq \max(\aleph_0, |F|)$.
By elementary set theory we can extend the given injection
$E \subset S$ to an injection $E' \to S$. In other words, we may
think of $E'$ as an element of our set $I$ contradicting the
maximality of $E$. This contradiction completes the proof.
\end{proof}

\begin{lemma}
\label{lemma-map-into-algebraic-closure}
Let $F$ be a field. Let $\overline{F}$ be an algebraic closure of $F$.
Let $M/F$ be an algebraic extension. Then there is a morphism of
$F$-extensions $M \to \overline{F}$.
\end{lemma}

\begin{proof}
Consider the set $I$ of pairs $(E, \varphi)$ where $F \subset E \subset M$
is a subextension and $\varphi : E \to \overline{F}$ is a morphism of
$F$-extensions. We partially order the set $I$ by declaring
$(E, \varphi) \leq (E', \varphi')$ if and only if $E \subset E'$ and
$\varphi'|_E = \varphi$. If $T = \{(E_t,  \varphi_t)\} \subset I$
is a totally ordered subset, then
$\bigcup \varphi_t : \bigcup E_t \to \overline{F}$ is an element of $I$.
Thus every totally ordered subset of $I$ has an upper bound.
By Zorn's lemma there exists a maximal element $(E, \varphi)$ in $I$.
We claim that $E = M$, which will finish the proof. If not, then
pick $\alpha \in M$, $\alpha \not \in E$. The $\alpha$ is algebraic
over $E$, see Lemma \ref{lemma-algebraic-goes-up}.
Let $P$ be the minimal polynomial of $\alpha$ over $E$.
Let $P^\varphi$ be the image of $P$ by $\varphi$ in $\overline{F}[x]$.
Since $\overline{F}$ is algebraically closed there is a root $\beta$
of $P^\varphi$ in $\overline{F}$. Then we can extend $\varphi$ to
$\varphi' : E(\alpha) = E[x]/(P) \to \overline{F}$ by mapping
$x$ to $\beta$. This contradicts the maximality of $(E, \varphi)$
as desired.
\end{proof}

\begin{lemma}
\label{lemma-algebraic-closures-isomorphic}
Any two algebraic closures of a field are isomorphic.
\end{lemma}

\begin{proof}
Let $F$ be a field. If $M$ and $\overline{F}$ are algebraic closures of
$F$, then there exists a morphism of $F$-extensions
$\varphi : M \to \overline{F}$ by
Lemma \ref{lemma-map-into-algebraic-closure}.
Now the image $\varphi(M)$ is algebraically closed.
On the other hand, the extension $\varphi(M) \subset \overline{F}$
is algebraic by Lemma \ref{lemma-algebraic-goes-up}.
Thus $\varphi(M) = \overline{F}$.
\end{proof}





